\documentclass[tikz,border=4pt]{standalone}
\usepackage{ctex}
\usepackage{amsmath}
\usetikzlibrary{calc, arrows.meta}

\begin{document}
% 三种情形：相交/平行/重合
\begin{tikzpicture}[scale=0.85, thick]

  %--- 情形1：相交（唯一解）---
  \begin{scope}[xshift=0cm]
    \draw[->] (-0.3,0) -- (3.2,0) node[right] {\small$x$};
    \draw[->] (0,-0.3) -- (0,3.2) node[above] {\small$y$};
    \node[below left, font=\tiny] at (0,0) {$O$};
    % 直线1
    \draw[blue, thick] (0.2,2.8) -- (2.8,0.2);
    % 直线2
    \draw[red, thick] (0.2,0.6) -- (2.4,2.8);
    % 交点
    \fill (1.4,1.6) circle (2pt);
    \node[below right, font=\small] at (1.5,-0.7) {相交：唯一解};
  \end{scope}

  %--- 情形2：平行（无解）---
  \begin{scope}[xshift=4.5cm]
    \draw[->] (-0.3,0) -- (3.2,0) node[right] {\small$x$};
    \draw[->] (0,-0.3) -- (0,3.2) node[above] {\small$y$};
    \node[below left, font=\tiny] at (0,0) {$O$};
    \draw[blue, thick] (0.5,2.8) -- (2.5,0.4);
    \draw[red, thick] (1.2,2.8) -- (3.2,0.4);
    \node[below right, font=\small] at (1.5,-0.7) {平行：无解};
  \end{scope}

  %--- 情形3：重合（无穷多解）---
  \begin{scope}[xshift=9.0cm]
    \draw[->] (-0.3,0) -- (3.2,0) node[right] {\small$x$};
    \draw[->] (0,-0.3) -- (0,3.2) node[above] {\small$y$};
    \node[below left, font=\tiny] at (0,0) {$O$};
    \draw[blue, thick, line width=3pt] (0.5,2.8) -- (2.5,0.4);
    \draw[red, dashed, thick] (0.5,2.8) -- (2.5,0.4);
    \node[below right, font=\small] at (1.5,-0.7) {重合：无穷多解};
  \end{scope}

\end{tikzpicture}
\end{document}
